\documentclass[a4paper]{article}
\usepackage{amsfonts}
\usepackage{amsmath}
\usepackage{amssymb}
\usepackage{graphicx}

\begin{document}

\noindent \large Auteur: Rob Willemsen \\
\noindent \large Studierichting: Informatica\\
\noindent \large Oefeningenreeks 1D nr. 19.2\\

\medskip

\normalsize

$2z^3 - z^2 + 2z - 1 = z^2(2z - 1) + 1(2z - 1)$ \\

$2z^3 - z^2 + 2z - 1 = (2z - 1)(z^2 + 1)$ \\

$2z - 1 = 0 \Rightarrow z = \dfrac{1}{2}$ \\

$z^2 + 1 = 0 \Rightarrow z^2 = -1 \Rightarrow z = \pm i$ \\

\textbf{Antwoord:} De nulpunten in $\mathbb{C}$ zijn $\dfrac{1}{2}, i$ en $-i$. De ontbinding over $\mathbb{R}[z]$ is $A(z) = (2z - 1)(z^2 + 1)$. \\

\begin{thebibliography}{999}
%\bibitem{a} Referentie
\end{thebibliography}
\end{document}